\section{High two-power strict vanishing at a growing range of coefficient primes}
\label{sec:effective-two-power-all-odd-primes}

\paragraph{The field used for Fourier coefficients remains fixed.}
For each coefficient prime $\ell$, the effective nonvanishing argument
uses $\mathbb F_{\ell^2}$, independently of the growing two-power
character.  The original coefficient ring of the strict group may
have arbitrarily large degree and ramification.  These two fields
are not identified.  The auxiliary cyclotomic class field below uses
a least common multiple: at $\ell=7$ it contains $\mu_7$, not
$\mu_{49}$.

Put $K=\mathbb Q(\sqrt{-3})$.  For each prime $\ell\geq5$ define
\[
 n_\ell=v_2(\ell^2-1),\qquad
 A_\ell=21\ell\,2^{2n_\ell},\qquad
 c_\ell=\min\{c\in\mathbb Z_{\geq0}:2^c>A_\ell\}.
 \tag{G175.1}
\]
Thus $n_\ell\geq3$ and $c_\ell\geq13$.
Write $c_{\rm cyc}$ for the $\ell$-adic cyclotomic character and
$\omega_\ell$ for its Teichmuller part; these are distinct from
the integer $c_\ell$ in (G175.1).
For a discrete $\ell$-primary coefficient $D$ set
\[
 \mathcal S_{\rm str}(K,D)=
 \ker\!\left(H^1_{\rm cts}(G_K,D)\longrightarrow
  \prod_{v\mid\ell}H^1(G_{K_v},D)
       \times\prod_{v\nmid\ell,\infty}H^1(I_v,D)\right).
 \tag{G175.2}
\]
The condition at every place above the coefficient prime is full
local vanishing.

\begin{theorem}[AC-G175: uniform coefficient-prime strict vanishing]
\label{AC-G175}
Let $\ell\geq5$ be prime, let $\delta$ be a primitive Dirichlet
character of exact conductor $2^c$ with $c\geq c_\ell$, and let
$\epsilon_7$ be any of the six characters modulo seven, including
the trivial primitive character.  Restrict these characters to $G_K$.
Let $E/\mathbb Q_\ell$ be any finite extension whose ring of integers
$O$ contains their values, and let $\mathfrak m_O$ be its maximal ideal.
For every $s\in\mathbb Z$, $e\geq1$, and continuous character
$\kappa:G_K\to1+\mathfrak m_O$ unramified at all finite places away
from $\ell$, one has
\[
 \begin{gathered}
 D=(E/O)(\delta\epsilon_7\omega_\ell^{-s}
                          \kappa^{-1}c_{\rm cyc}^{-e}),\\
                         \mathcal S_{\rm str}(K,D)=0.
 \end{gathered}
 \tag{G175.3}
\]
The residual strict group is zero over every finite residue field
containing its character values.  No residue of $e+s$ modulo
$\ell-1$ is excluded.  The stronger, simpler numerical condition
\[
                         2^c>84\ell(\ell+1)^2
 \tag{G175.4}
\]
is sufficient for (G175.3).

In the actual host application take $s=0$ and $\kappa=1$.
The finite host label is exactly $\delta\epsilon_7$, with trivial
cubic and auxiliary quadratic factors.  No finite conductor-$\ell$
character is inserted into that label, and no original coefficient
ring or lattice is replaced.
\end{theorem}

\begin{proof}
\proofstep{1. Effective special-value nonvanishing over $\mathbb F_{\ell^2}$.}
The effective nonvanishing theorem (G173.1)--(G173.2) is stated for
every odd coefficient prime.  Its complete proof uses Sinnott's
Lemma~3.6 and Proposition~4.6, including the one-character trace
argument and avoidance of the exceptional Fourier roots.
We apply that theorem at $\ell$, not its later specialized strict
conclusion at $269$.

Since $c\geq13$, the unique decomposition
\[
          \delta=\chi_{-4}^{\,a}\psi,
       \qquad a\in\{0,1\},\quad\psi(-1)=1
 \tag{G175.5}
\]
has $\psi$ primitive of the same exact conductor $2^c$.
Every fixed factor needed for the class calculation is of the form
\[
        \theta=\chi_{-4}^{\,a}\chi_{-3}^{\,b}
                    \epsilon_7^{\,\pm1}\omega_\ell^j,
          \qquad b\in\{0,1\},\ j\bmod\ell-1.
 \tag{G175.6}
\]
We use only the odd factors.  Their primitive conductors divide
$\operatorname{lcm}(84,\ell)$.  This statement uses a least common
multiple also when $\ell=7$; no square of seven is required.
Their reductions, and fourth roots of unity, belong to
$k=\mathbb F_{\ell^2}$: $\ell^2-1$ is divisible by $4$ and $6$,
and the Teichmuller values belong to $\mathbb F_\ell$.
The field $k$ is fixed before $c$ and the values of $\psi$ are chosen.

If $f_\theta$ is the primitive conductor, then
\[
 L_\theta=\operatorname{lcm}(f_\theta,2^{n_\ell})
       \leq21\ell\,2^{n_\ell},\qquad
                         v_2(4f_\theta)\leq4.
 \tag{G175.7}
\]
The definition of $c_\ell$ implies
$c\geq2n_\ell$ and $c>4$, and also
\[
              2^{c-n_\ell}>21\ell\,2^{n_\ell}
                                     \geq L_\theta>L_\theta-1.
 \tag{G175.8}
\]
Thus all hypotheses of (G173.2) hold.  It proves that every odd
character consisting of (G175.6) times $\psi$ or $\psi^{-1}$ has
$\ell$-unit value at zero.  This includes integrality, even if the
primitive conductor contains $\ell$.  It holds under each embedding;
subsequent extension of residue fields preserves the unit property.

For the simpler sufficient bound, one of $\ell-1,\ell+1$ has
two-adic valuation one.  If the other's valuation is $a$, then
$n_\ell=a+1$ and $2^a\leq\ell+1$.  Consequently
\[
 2^{n_\ell}\leq2(\ell+1),\qquad
                   A_\ell\leq84\ell(\ell+1)^2.
 \tag{G175.9}
\]
This proves the implication from (G175.4).

\proofstep{2. Primitive integral annihilators for the odd class sectors.}
Set
\[
 \begin{gathered}
 m_{\ell,c}=\operatorname{lcm}(21\cdot2^c,\ell),\quad
 H=\mathbb Q(\zeta_{m_{\ell,c}}),\\
 G=\operatorname{Gal}(H/\mathbb Q),\quad
 J=\operatorname{Gal}(H/K),\quad A_H=\operatorname{Cl}(H)\{\ell\}.
 \end{gathered}
 \tag{G175.10}
\]
The degree of $H$ is
\[
 [H:\mathbb Q]=
 \begin{cases}
  12\cdot2^{c-1},&\ell=7,\\
  12\cdot2^{c-1}(\ell-1),&\ell\ne7.
 \end{cases}
 \tag{G175.11}
\]
It is prime to $\ell$ in both cases.  In particular integral
character idempotents exist after a finite unramified coefficient
extension containing the values of all characters of $G$.

Let $\nu$ be an odd product of a fixed factor in (G175.6) and
an even primitive character of conductor $2^c$.
At $\ell=7$, the $\epsilon_7$ and $\omega_7$ factors are evaluated
on the same conductor-seven quotient; they are not independent
cyclotomic factors.
Its primitive conductor $f_\nu$ divides $m_{\ell,c}$.
Work first over $H_\nu=\mathbb Q(\zeta_{f_\nu})$.
Choose a positive integer $b_0$ by the Chinese remainder theorem with
\[
 b_0\equiv5\pmod {2^c},\qquad
                  b_0\equiv1\pmod {\operatorname{lcm}(21,\ell)}.
 \tag{G175.12}
\]
Then $b_0$ is prime to $f_\nu$ and $b_0\equiv1\pmod\ell$.
The integral Stickelberger element
\[
 (b_0-\sigma_{b_0})\Theta_{f_\nu},\qquad
 \Theta_f=f^{-1}\sum_{\substack{1\leq a\leq f\\(a,f)=1}}
                                      a\sigma_a^{-1}
 \tag{G175.13}
\]
annihilates the class group of $H_\nu$.
This is the literal inverse-group-element convention in Sharifi's
Definition~4.5.1, Lemma~4.5.2 and Theorem~4.5.6.
Its scalar on the $\nu$ component is
\[
             -(b_0-\nu(b_0))L(0,\nu^{-1}).
 \tag{G175.14}
\]
The second factor is a unit by Step~1.  The first has reduction
$1-\overline\psi(5)$, with $\psi$ replaced by its inverse if
appropriate.  This is nonzero: $5$ generates
$1+4\mathbb Z/2^c\mathbb Z$, so $\psi(5)$ has exact order
$2^{c-2}$, and reduction in odd characteristic preserves that order.
Thus the $\nu$ class component at primitive conductor is zero.
We never apply an imprimitive Bernoulli scalar and discard its Euler
factors, nor use the isolated rational element $\Theta_f$ integrally.

For $V=\operatorname{Gal}(H/H_\nu)$, ideal extension $j$ and norm
$N$ satisfy $Nj=|V|$ and $jN=\sum_{v\in V}v$.
The order of $V$ is prime to $\ell$, and $\nu$ is trivial on $V$.
These identities identify the $\nu$ component at the two levels,
so it is zero in $A_H$ as well.  This is exactly the applied
norm/idempotent argument of Sharifi's Lemma~1.3.5.

\proofstep{3. Reflection in the actual least-common-multiple field.}
The field $H$ is CM, abelian over $\mathbb Q$, contains $\mu_\ell$,
and has degree prime to $\ell$.
Let $B\subset H^\times$ be the elements whose $\ell$th roots
generate everywhere unramified extensions of $H$.
The equivariant Kummer/class-field construction gives an exact
sequence
\[
 0\longrightarrow
 (B\cap\mathcal O_H^\times)/\mathcal O_H^{\times\ell}
 \longrightarrow\operatorname{Hom}(A_H/\ell A_H,\mu_\ell)
 \longrightarrow A_H[\ell].
 \tag{G175.15}
\]
No surjectivity at the last term is asserted.
This is the sequence in Sharifi's Lemma~1.4.6 and the complete
proof of Theorem~4.4.2, with the Artin equivariance of
Proposition~1.3.2.

In a CM field the odd part of units modulo torsion is zero, by
the unit theorem; the only possible odd line modulo $\ell$ is
the $\ell$-primary root-of-unity line of character $\omega_\ell$.
Here the $\ell$-primary roots in $H$ are precisely $\mu_\ell$.
At every place $v\mid\ell$, the conductor other than $\ell$ is
unramified, whereas the single factor $\ell$ gives
\[
                    e(H_v/\mathbb Q_\ell)=\ell-1.
 \tag{G175.16}
\]
Adjoining $\mu_{\ell^2}$ has relative ramification index $\ell$,
unchanged by the unramified base extension.  The root-of-unity
line therefore cannot belong to the unramified radical kernel
in (G175.15).  This proves its entire odd kernel is zero,
including its $\omega_\ell$ component.
At $\ell=7$ this argument would fail for the erroneous modulus
$21\cdot2^c\ell$, which inserts $49$ and destroys (G175.11);
the modulus in (G175.10) is essential.

For an even character $\nu$, the part of the middle term dual
to $A_H(\nu)$ has character $\omega_\ell\nu^{-1}$.
It injects into $A_H[\ell](\omega_\ell\nu^{-1})$.  Thus
\[
 A_H(\omega_\ell\nu^{-1})=0\quad\Longrightarrow\quad A_H(\nu)=0.
 \tag{G175.17}
\]
The reflected character is odd, still has exact two-primary
conductor $2^c$, and is again of the form just treated, with
$\epsilon_7$ inverted and the Teichmuller exponent changed.
Step~2 kills it.  We have therefore killed every required even
and odd component.  We have not asserted that the whole class
group $A_H$ is zero.

\proofstep{4. Residual strict classes use the original, nondual character.}
Let $k'$ be a finite residue field containing the character values,
and write
\[
 \alpha=\overline\delta\,\overline\epsilon_7\,
                \overline\omega_\ell^{-r},
                  \qquad r\equiv e+s\pmod{\ell-1}.
 \tag{G175.18}
\]
Restriction from $G_K$ to $G_H$ is injective in degree one,
because $[H:K]$ is prime to $\ell$.
Over $H$ the coefficient is trivial.  If $z$ is the restricted
cocycle, then for $g\in G_H$ and $\tau\in G_K$,
\[
                   z(\tau g\tau^{-1})=\alpha(\tau)z(g).
 \tag{G175.19}
\]
Every strict class vanishes on all inertia and on the full
decomposition groups above $\ell$.  Its homomorphism therefore
factors through a quotient of the ordinary class group, and gives
\[
 \mathcal S_{\rm str}(K,k'(\alpha))\lhook\joinrel\longrightarrow
 \operatorname{Hom}_{\mathbb F_\ell[J]}
                          (A_H/\ell A_H,k'(\alpha)).
 \tag{G175.20}
\]
The character selected here is $\alpha$, not $\alpha^{-1}$.
The two rational extensions of $\alpha$ differ by $\chi_{-3}$;
both are among the components killed in Steps~2--3.
After extension to a residue field containing the character
idempotents, the right side is zero.  Faithful scalar extension
then proves it was zero over the original $k'$.
This proves residual strict zero for every residue $r$, without
identifying $k'$ with the fixed field $\mathbb F_{\ell^2}$.

\proofstep{5. High two-power Frobenius excludes every local invariant.}
Since $\ell\ne3$, every $K_v/\mathbb Q_\ell$ for $v\mid\ell$
is unramified of degree $f_v\in\{1,2\}$.
Write $\ell=\tau u$ with $\tau\in\{1,-1\}$ and
$u\in1+4\mathbb Z_2$.  Then
\[
 a=v_2(u-1)\geq2,\qquad n_\ell=a+1.
 \tag{G175.21}
\]
For a primitive even $\psi$ of conductor $2^c$, the order of
$\psi(u)$ is $2^{c-a}$: writing $u=5^z$ gives
$v_2(z)=a-2$, and $\psi(5)$ has order $2^{c-2}$.
The possible parity sign in $\delta(\ell)$ cannot alter this
order when it is greater than two.  Our bound on $c$ gives
\[
 \operatorname{ord}\delta(\ell)=2^{c-n_\ell+1},\qquad
 \operatorname{ord}\delta(\ell)^{f_v}
                            \geq2^{c-n_\ell}>2.
 \tag{G175.22}
\]

Choose an arithmetic Frobenius lift in $G_{K_v}$ that fixes
$\mu_\ell$.  Such a lift exists because $K_v/\mathbb Q_\ell$
is unramified and $\mathbb Q_\ell(\mu_\ell)/\mathbb Q_\ell$
is totally ramified.  If $\ell\ne7$, the finite factor on it is
\[
                        \delta(\ell)^{f_v}\epsilon_7(\ell)^{f_v}.
 \tag{G175.23}
\]
The second factor has order dividing six, with two-primary part
of order at most two.  It cannot cancel the first factor in
(G175.22).  If $\ell=7$, the chosen lift fixes $\mu_7$, so
$\epsilon_7$ is exactly one on it; the first factor is still
nontrivial.  All these finite values have order prime to $\ell$,
so their reductions retain the relevant orders.

On this same lift the residual cyclotomic and Teichmuller factors
are one, and $\kappa$ reduces to one.  The scalar of the actual
coefficient $D$ minus one is therefore an $O$-unit.  This proves
\[
                         D^{G_{K_v}}=0\qquad(v\mid\ell)
 \tag{G175.24}
\]
for every residue $r$, including zero.  Global invariants vanish
as well.  Inverting the Frobenius convention inverts the same
nontrivial finite scalar and changes none of this conclusion.

At $v\nmid\ell$, all cyclotomic, Teichmuller and $\kappa$ factors
are unramified.  The remaining inertia action has finite order
dividing a power of two times six, hence prime to $\ell$.
Its invariants on $D$ are either all of $D$ or zero, and in
either case are divisible by an original uniformizer $\varpi$.

\proofstep{6. Lift using the original uniformizer.}
The module $D$ is $\varpi$-divisible and
$D[\varpi]=k'(\alpha)$.
Apply cohomology to
$0\to D[\varpi]\to D\xrightarrow{\varpi}D\to0$
globally, at each $G_{K_v}$ above $\ell$, and at each other
finite inertia group.  The invariant calculations in Step~5
make the local maps from residual $H^1$ injective.
Thus every strict class killed by $\varpi$ lifts to a strict
residual class.  Step~4 gives
$\mathcal S_{\rm str}(K,D)[\varpi]=0$.
Every continuous cocycle from a profinite group to the discrete
module $D$ has finite image and is killed by a power of
$\varpi$.  Hence the whole strict group is zero, as claimed.
No absolute-ramification restriction on $O$ was used.
\end{proof}

\paragraph{Exact actual-host scope and the finite growing prime range.}
For the actual host in (FG.3), its two-primary modulus is
$2^{h+2}$ and its group exponent over $K$ is $N=3\cdot2^h$.
Whenever $c_\ell\leq c\leq h+2$, every character
$\delta\epsilon_7$ in the theorem is an actual label with
trivial cubic and auxiliary quadratic factors.
Its order divides $N$, so its values belong to the original
$O_\ell$, the prescribed completion of $\mathbb Z[\zeta_N]$.
The coefficient in (FC.2) is exactly (G175.3) with $s=0$ and
$\kappa=1$.  Inflation from its prescribed ramification set into
$G_K$ is injective and preserves all the displayed local zeros;
outside the set the inflated cocycle already vanishes on inertia.
Thus its original strict group is zero.

For a fixed $h$, the exact number of such labels at a given
$\ell\geq5$ is
\[
               C_\ell(h)=6\max\{0,2^{h+1}-2^{c_\ell-2}\}.
 \tag{G175.25}
\]
Indeed there are $2^{c-2}$ primitive characters of conductor
$2^c$ and six choices at seven, and their restrictions to $K$
are distinct.  Only finitely many coefficient primes occur:
$n_\ell\geq3$ and $c_\ell\leq h+2$ imply
$1344\ell<2^{h+2}$.
The simpler condition $84\ell(\ell+1)^2<2^{h+2}$ describes
a uniform growing subrange without any valuation of $\ell\pm1$.
For example $n_5=3$ and $c_5=13$, so the new coefficient-five
family is nonempty for $h\geq11$.

These are pointwise strict statements, independent of the prescribed
pair $q,d$ and every weight cutoff.  They do not make the derived
indices individually zero.  Their contribution to an actual
omitted-term union must still be compared with (EU.2) and the one
full-union (FL.4) bound on the identical window.
No theorem for coefficients two or three, mixed-cubic labels,
growing auxiliary quadratic twists, or the entire $I_{399}$
remainder is asserted here.

\paragraph{Source and geometric scope.}
The applied primary reflection and annihilator proofs are the full
passages Sharifi~1.3.1--1.3.5, 1.4.4--1.4.6,
1.4.12--1.4.15, 4.4.2 and 4.5.1--4.5.6.
Class-field theory, the unit theorem and the stated foundational
inputs to those printed proofs remain imported.
The new coefficient-prime argument rederives the least-common-multiple
field, local ramification, fixed finite field and original-uniformizer
steps above; it does not silently extend the specialized theorem at269.
The rank-one finite characters and unramified Kummer radicals use the
previously audited finite-etale and $\mu_\ell$-torsor dictionary.
No new motive, regulator comparison, analytic realization or
Lean execution is claimed.
